% !TEX root = task-2-2-2.tex
% Editable task bank. Append new BankTask blocks at the end of the document.
\documentclass[12pt,a4paper]{article}
\usepackage{physics-textbook}
\usepackage{graphicx,needspace}
\usetikzlibrary{decorations.pathmorphing}
\newif\ifShowSolutions
\ifdefined\TutorEdition
  \ShowSolutionstrue
\else
  \ShowSolutionsfalse % Change to \ShowSolutionstrue to display solutions.
\fi
\newcounter{banktask}
\newsavebox{\BankTaskBox}
\newcommand{\BankTaskContent}[3]{%
  \noindent\textbf{\thebanktask.}\quad #2\par
  \ifShowSolutions
    \Needspace{50mm}% Keep the solution heading with its opening text or diagram.
    #3\par
    \ifdefined\TutorEdition
      \ifcsname TutorSource#1\endcsname
        \nopagebreak\smallskip{\small\csname TutorSource#1\endcsname\par}%
      \fi
    \fi
  \fi
}
\newcommand{\BankTask}[3]{%
  \refstepcounter{banktask}\par\medskip
  \sbox{\BankTaskBox}{\begin{minipage}{\linewidth}%
    \BankTaskContent{#1}{#2}{#3}\end{minipage}}%
  \ifdim\dimexpr\ht\BankTaskBox+\dp\BankTaskBox\relax>\dimexpr\textheight-20mm\relax
    \Needspace{35mm}\BankTaskContent{#1}{#2}{#3}%
  \else
    \Needspace{\dimexpr\ht\BankTaskBox+\dp\BankTaskBox\relax}%
    \noindent\usebox{\BankTaskBox}\par
  \fi
  \medskip
}
\begin{document}
\BookSubsection{2.2.2}{Твёрдые тела}


% Source sdamgia-810: https://phys-ege.sdamgia.ru/problem?id=810
\BankTask{sdamgia-810}{
Какое из утверждений справедливо для кристаллических тел?

$1$) во время плавления температура кристалла изменяется

$2$) в расположении атомов кристалла отсутствует порядок

$3$) атомы кристалла расположены упорядоченно

$4$) атомы свободно перемещаются в пределах кристалла
}{
\smallskip\noindent\textbf{Решение}\par
Кристалл — это агрегатное состояние вещества, которое характеризуется строгим порядком расположения атомов. Атомы в кристаллах не могут свободно перемещаться, колеблясь около узлов кристаллической решетки. Во время плавления температура кристалла не меняется, вся поступающая энергия идет на разрушение кристаллической решетки. Верно утверждение $3$.

Ответ: $3$.
}

% Source sdamgia-828: https://phys-ege.sdamgia.ru/problem?id=828
\BankTask{sdamgia-828}{
В процессе перехода вещества из жидкого состояния в кристаллическое

$1$) существенно увеличивается расстояние между его молекулами

$2$) молекулы начинают притягиваться друг к другу

$3$) существенно увеличивается упорядоченность в расположении его молекул

$4$) существенно уменьшается расстояние между его молекулами
}{
\smallskip\noindent\textbf{Решение}\par
В процессе перехода вещества из жидкого состояния в твердое происходит построение кристаллической решетки, а значит, существенно увеличивается упорядоченность в расположении молекул вещества.

Ответ: $3$.
}

% Source sdamgia-836: https://phys-ege.sdamgia.ru/problem?id=836
\BankTask{sdamgia-836}{
Какие частицы находятся в узлах решетки металла?

$1$) нейтральные атомы

$2$) электроны

$3$) отрицательные частицы

$4$) положительные ионы
}{
\smallskip\noindent\textbf{Решение}\par
В узлах решетки металла находятся положительные ионы. Валентные электроны в атомах металлов слабо связаны с ядром. Поэтому в металле они обобществляются, перестают принадлежать определенному атому и образуют так называемый «электронный газ». Атомы металла в узлах решетки, потерявшие свои валентные электроны, превращаются в положительные ионы.

Ответ: $4$.
}

% Source sdamgia-837: https://phys-ege.sdamgia.ru/problem?id=837
\BankTask{sdamgia-837}{
Какое свойство отличает монокристалл от аморфного тела?

$1$) прочность

$2$) электропроводность

$3$) прозрачность

$4$) анизотропность
}{
\smallskip\noindent\textbf{Решение}\par
Аморфное тело не имеет кристаллической структуры и, как правило, — изотропно, то есть не обнаруживают различных свойств в разных направлениях. Монокристаллом же называется отдельный однородный кристалл, имеющий непрерывную кристаллическую решетку и характеризующийся анизотропией свойств. Таким образом, главным свойством, отличающим монокристалл от аморфного тела, является анизотропность.

Ответ: $4$.
}

% Source sdamgia-3789: https://phys-ege.sdamgia.ru/problem?id=3789
\BankTask{sdamgia-3789}{
Какое из приведенных ниже утверждений справедливо для кристаллических тел?

$1$) в расположении атомов отсутствует порядок

$2$) атомы свободно перемещаются в пределах тела

$3$) при изобарном плавлении температура тела остается постоянной

$4$) при одинаковой температуре диффузия в кристаллах протекает быстрее, чем в газах
}{
\smallskip\noindent\textbf{Решение}\par
Твердые тела обладают кристаллической решеткой, в расположении атомов присутствует дальний порядок. Атомы располагаются в узлах решетки, и в отличие от газов не могут свободно перемещаться по объему тела. Это в частности приводит к тому, что при прочих равных условиях диффузия в твердых телах протекает гораздо медленнее, чем в газах. Таким образом, утверждения $1$, $2$ и $4$ не верны. А вот утверждение $3$ справедливо. Опыт показывает, что при изобарном плавлении, температура твердого тела не изменяется, все поступающее к телу тепло идет на разрушение кристаллической решетки.

Ответ: $3$.
}

% Source sdamgia-4343: https://phys-ege.sdamgia.ru/problem?id=4343
\BankTask{sdamgia-4343}{
Модель, служащая для демонстрации внутреннего строения тел, устроена следующим образом: в дно прямоугольной коробки воткнуты одинаковые упругие вертикальные стерженьки, на каждый из которых насажен магнитик в виде плоской таблетки. После приведения одного из магнитиков в колебательное движение вскоре начинают хаотически колебаться на стерженьках и все остальные магнитики, отталкиваясь друг от друга. Эта модель лучше всего иллюстрирует поведение молекул

$1$) идеального газа

$2$) жидкости

$3$) твердого тела

$4$) идеального газа и жидкости
}{
\smallskip\noindent\textbf{Решение}\par
Данная модель лучше всего демонстрирует поведение молекул твердого тела, потому что именно в данном агрегатном состоянии частицы совершают малые колебания вокруг положений равновесия, не перемещаясь по объему тела.

Ответ: $3$
}

% Source sdamgia-7695: https://phys-ege.sdamgia.ru/problem?id=7695
\BankTask{sdamgia-7695}{
В стакане находится вода в кристаллическом состоянии при температуре $0^\circ$C. Из-за контакта с теплым стаканом и с комнатным воздухом вся вода переходит в жидкое состояние. В процессе перехода воды из кристаллического состояния в жидкое

$1$) ее объем уменьшается;

$2$) ее плотность уменьшается;

$3$) ее внутренняя энергия уменьшается;

$4$) ее температура уменьшается.
}{
\smallskip\noindent\textbf{Решение}\par
При переходе воды из кристаллического состояния в жидкое, ее плотность увеличивается, а значит при неизменной массе вещества объем должен уменьшиться.

Ответ: $1$
}


% Source sdamgia-4155: https://phys-ege.sdamgia.ru/problem?id=4155
\BankTask{sdamgia-4155}{
Аморфными называют тела, 

$1$) атомная структура которых имеет ближний порядок и не имеет дальнего порядка

$2$) которые находятся в твердом состоянии и представляют собой один большой кристалл

$3$) которые находятся в твердом состоянии и состоят из большого числа хаотически расположенных мелких кристаллов

$4$) которые находятся в жидком состоянии
}{
\smallskip\noindent\textbf{Решение}\par
Аморфными называют тела, атомная структура которых имеет ближний порядок и не имеет дальнего порядка, характерного для кристаллических структур.

Ответ: $1$.
}

% Source sdamgia-5465: https://phys-ege.sdamgia.ru/problem?id=5465
\BankTask{sdamgia-5465}{
Молекулы любого твердого тела находятся в непрерывном движении. Почему тогда тела не рассыпаются на отдельные частицы?

$1$) нет сил отталкивания между молекулами

$2$) среднее значение скорости молекул равно нулю

$3$) действует взаимное притяжение молекул

$4$) действует атмосферное давление
}{
\smallskip\noindent\textbf{Решение}\par
Молекулы твердого тела колеблются в узлах кристаллической решетки вокруг положения равновесия. В этом положении взаимное притяжение молекул сильнее, чем силы отталкивания.

Ответ: $3$.
}

% Source sdamgia-4366: https://phys-ege.sdamgia.ru/problem?id=4366
\BankTask{sdamgia-4366}{
Зимой школьник решил поставить опыт: полностью заполнил две тонкие пластиковые бутылки с практически нерастяжимыми стенками горячей водой (почти кипятком), потом из одной вылил воду, сразу же обе плотно закрыл крышками и выставил бутылки на мороз на всю ночь. В результате одна бутылка лопнула, а другая сплющилась. Объясните, основываясь на известных физических законах и закономерностях, какая из бутылок лопнула и почему.
}{
\smallskip\noindent\textbf{Решение}\par
$1$. В первой бутылке были насыщенные пары воды и малое количество воздуха, а во второй — только горячая вода.

$2$. Ночью на морозе обе бутылки за счет всех видов теплоотдачи — конвекции, теплопроводности и теплового излучения — остыли до температуры окружающего холодного воздуха.

$3$. Пары воды в первой бутылке сконденсировались и замерзли, сильно уменьшив свой объем, а давление в бутылке упало до очень низкого значения. Поэтому бутылка под действием наружного атмосферного давления сплющилась.

$4$. Во второй бутылке вода сначала охладилась до $0^\circ$C, а потом при фазовом переходе из жидкого состояния в твердое отдала теплоту кристаллизации, замерзла и охладилась до низкой температуры окружающей среды. При этом объем воды вырос примерно на $10$\% (это следует из таблицы плотностей веществ), что привело к сильному росту давления в бутылке с нерастяжимыми стенками, и она лопнула.
}

\end{document}
